\begin{abstract}
Rapid episodic memory requires not only selecting stored content, but expressing
that content in a form that downstream circuits can interpret. We examine this
constraint in a minimal hippocampal model in which entorhinal cortex (EC)
generates a CA3-like retrieval key and an instructed CA1 state, plastic
CA3--CA1 weights associate the key with that state, and a fixed pretrained
CA1--EC decoder reconstructs the input. Pretraining establishes a
representational basis before rapid associative learning. Across 20 paired
initializations, a fixed permutation of the instructed CA1 coordinates impaired
one-shot reconstruction under both an instructive-driven rule (mean
output--target cosine, 0.804 aligned versus 0.125 permuted) and a bounded
error-driven rule (0.514 versus 0.066). Applying the corresponding permutation
to the decoder restored performance, demonstrating that preserved information
can become unreadable to an unchanged output pathway. On a cue-guided circular
track, exchanging cue identities caused abrupt, reproducible changes in CA1
tuning relative to a matched no-swap control. Event-aligned and population-map
analyses revealed a transient update at exchange boundaries and alternating
representational blocks tied to cue assignment. Both rules generated
continuous mixtures of spatially stable and cue-modulated responses, while differing
quantitatively in one-shot reconstruction, tuning stability, and position
decodability. During frozen recall, selective degradation of lateral or medial
EC inputs preferentially impaired cue or position readout, respectively, while
the complementary cue readout was preserved under medial EC degradation. A
dense common-key control failed to distinguish cue or position states. These
simulations support a limited proposal: CA3-like keys can select associative
content, while rapid CA1 plasticity expresses that content within a structured
basis that remains legible to an established readout. The model demonstrates
degraded-cue robustness, not recurrent CA3 pattern completion or long-term
stability.
\end{abstract}
